% =====================================================================
%  2bat-ud3-9.tex  —  SESSIÓ 9: On fa un salt la prestació? (estudi de cas)
% =====================================================================
\horatitol{Sessió 9 — On fa un salt la prestació?}%
{Estudi de cas (Activitat 3): comprovar la continuïtat d'una prestació per trams i interpretar socialment el salt (l'efecte «esglaó»).}

\subsection{Idea clau}

Apliquem la continuïtat a un cas real: una \textbf{prestació per trams}, com les
funcions a trossos que vau veure a 1r. La pregunta clau: a les fronteres entre trams,
la prestació \textbf{enllaça suaument} o fa un \textbf{salt}? La resposta té una
lectura social directa.

\begin{idea}
\textbf{Comprovar la continuïtat a una frontera.} En el valor frontera, calculem la
prestació \textbf{des d'un tram i des de l'altre}:
\begin{itemize}[leftmargin=1.4em, itemsep=2pt]
  \item Si \textbf{coincideixen}, la funció és \textbf{contínua}: el canvi és suau.
  \item Si \textbf{no coincideixen}, hi ha un \textbf{salt}: discontinuïtat.
\end{itemize}
\textbf{Lectura social del salt (l'efecte «esglaó»).} Si hi ha un salt cap avall, una
persona que guanya \textbf{només una mica més} que el llindar pot perdre
\textbf{bruscament} una part important de l'ajut. És la \textbf{trampa de la pobresa}:
desincentiva millorar els ingressos. Una prestació \textbf{contínua} reparteix el
canvi de manera \textbf{suau i justa}.
\end{idea}

\subsection{Exemples}

\begin{exemple}[una prestació contínua (disseny just)]
\[
A(x)=
\begin{cases}
600 & \text{si } x<800,\\[2pt]
600-0,5\,(x-800) & \text{si } 800\le x<2000,\\[2pt]
0 & \text{si } x\ge 2000,
\end{cases}
\]
on $x$ són els ingressos mensuals. Comprovem les fronteres: a $x=800$, els dos trossos
donen $600$; a $x=2000$, el tram del mig dóna $600-0,5\cdot 1200=0$, igual que el
tercer. \textbf{Tot enllaça}: la prestació és \textbf{contínua}, decreix de manera
gradual.

\begin{center}
\begin{tikzpicture}
\begin{axis}[udaxis, width=9.4cm, height=4.4cm,
    xlabel={\small ingressos ($x$)}, ylabel={\small ajut (\euro)},
    xmin=0, xmax=2600, ymin=0, ymax=750, xtick={0,800,2000}, ytick={0,300,600}]
  \addplot[udblue, domain=0:800, samples=2]{600};
  \addplot[udblue, domain=800:2000, samples=2]{600-0.5*(x-800)};
  \addplot[udblue, domain=2000:2600, samples=2]{0};
  \addplot[udnavy, only marks, mark=*, mark size=1.4pt] coordinates {(800,600)(2000,0)};
\end{axis}
\end{tikzpicture}

\figcap{Prestació contínua: l'ajut baixa de mica en mica, sense salts.}
\end{center}
\end{exemple}

\begin{exemple}[una prestació amb salt (efecte esglaó)]
\[
B(x)=
\begin{cases}
600 & \text{si } x<\num{1000},\\[2pt]
0 & \text{si } x\ge \num{1000}.
\end{cases}
\]
A $x=\num{1000}$: pel costat esquerre val $600\eur$ i pel dret, $0\eur$. \textbf{Salt}
de $600$ a $0$: discontinuïtat. Una persona que passa de $999\eur$ a $\num{1000}\eur$
perd de cop \textbf{tot} l'ajut.

\begin{center}
\begin{tikzpicture}
\begin{axis}[udaxis, width=9.4cm, height=4.2cm,
    xlabel={\small ingressos ($x$)}, ylabel={\small ajut (\euro)},
    xmin=0, xmax=2000, ymin=0, ymax=750, xtick={0,1000}, ytick={0,300,600}]
  \addplot[udblue, domain=0:1000, samples=2]{600};
  \addplot[udblue, domain=1000:2000, samples=2]{0};
  \addplot[udnavy, only marks, mark=*, mark size=1.5pt] coordinates {(1000,0)};
  \addplot[udnavy, only marks, mark=o, mark size=1.5pt] coordinates {(1000,600)};
  \draw[udgrey, dashed] (axis cs:1000,0) -- (axis cs:1000,600);
  \node[udnavy, font=\scriptsize, anchor=west] at (axis cs:1040,300) {salt $600\eur$};
\end{axis}
\end{tikzpicture}

\figcap{Prestació amb salt: per $1\eur$ més, es perd tot l'ajut (efecte «esglaó»).}
\end{center}
\end{exemple}

\subsection{Exercicis resolts}

\begin{exresolt}[comprovar la continuïtat a una frontera]
Comprova si aquesta prestació és contínua a $x=\num{1000}$:
\[
A(x)=
\begin{cases}
500 & \text{si } x<\num{1000},\\[2pt]
500-0,4\,(x-\num{1000}) & \text{si } x\ge \num{1000}.
\end{cases}
\]
\solucio
A $x=\num{1000}$: pel costat esquerre, $A=500$; pel dret,
$500-0,4\cdot 0=500$. Tots dos donen $500$: \textbf{coincideixen}, la funció és
\textbf{contínua} a $x=\num{1000}$ (l'ajut comença a baixar suaument a partir d'aquí,
sense salt).
\resposta{Contínua: tots dos trossos valen $500$ a $x=\num{1000}$.}
\end{exresolt}

\begin{exresolt}[redactar una conclusió social]
Compara les prestacions $A$ (contínua) i $B$ (amb salt) dels exemples. Quina és més
justa i per què?
\solucio
La prestació \textbf{$A$ és més justa}: l'ajut decreix de manera gradual a mesura que
pugen els ingressos, de manera que guanyar una mica més \emph{sempre} deixa la persona
en una situació millor o igual. La prestació \textbf{$B$ no ho és}: el salt fa que, en
superar el llindar per $1\eur$, es perdi tot l'ajut de cop (efecte esglaó / trampa de
la pobresa), cosa que penalitza millorar els ingressos. Un bon disseny de prestacions
evita els salts.
\resposta{$A$, perquè és contínua i el canvi és gradual; $B$ té un salt injust que
penalitza superar el llindar per poc.}
\end{exresolt}

\subsection{Exercicis per fer}

\begin{exercicis}
\begin{exlist}
  \item Comprova si cada prestació és contínua a la frontera indicada:
        \begin{enumerate}[label=\alph*), leftmargin=1.6em, topsep=2pt]
          \item $A(x)=\begin{cases}400 & x<\num{1000}\\ 400-0,2(x-\num{1000}) & x\ge\num{1000}\end{cases}$ a $x=\num{1000}$
          \item $B(x)=\begin{cases}300 & x<\num{1500}\\ 100 & x\ge\num{1500}\end{cases}$ a $x=\num{1500}$
        \end{enumerate}
  \item Per a la prestació (b) anterior, quina és l'altura del salt? Què hi perd qui
        passa de $\num{1499}\eur$ a $\num{1500}\eur$?
  \item La prestació $A(x)=\begin{cases}500 & x<600\\ 350 & 600\le x<1200\\ 0 & x\ge 1200\end{cases}$
        és contínua? Localitza tots els salts i digues la seva altura.
  \item Dibuixa a mà una prestació per trams que sigui \textbf{contínua} (sense salts)
        i que vagi disminuint a mesura que pugen els ingressos.
  \item \textbf{(Conclusió social)} Per a la prestació de l'exercici 3, redacta dues o
        tres frases dient si et sembla justa i per què, fent servir la idea d'efecte
        esglaó.
  \item \textbf{(Disseny)} Proposa com modificaries la prestació $B$ de l'exercici 1
        (la que té un salt) perquè fos contínua, mantenint que l'ajut acabi a $0$ per a
        ingressos alts.
\end{exlist}
\end{exercicis}

% ---------------------------------------------------------------------
\fullsolucions{Sessió 9}
\begin{solucions}
\begin{sollist}
  \item (a) esquerre $400$, dret $400-0=400$: \textbf{contínua}. \quad
        (b) esquerre $300$, dret $100$: \textbf{discontínua}, salt de $300$ a $100$.
  \item Altura del salt: $300-100=200\eur$. Qui passa de $\num{1499}$ a $\num{1500}\eur$
        perd $200\eur$ d'ajut de cop (per $1\eur$ més d'ingressos).
  \item No és contínua: salt a $x=600$ (de $500$ a $350$, altura $150$) i a $x=1200$
        (de $350$ a $0$, altura $350$).
  \item Esbós: una funció esglaonada \emph{sense} salts, per exemple constant i després
        un tram que baixa en línia recta fins a $0$, enganxant-se a les fronteres.
  \item Resposta oberta. Ha d'identificar que té salts (a $600$ i $1200$) i argumentar
        que penalitza superar els llindars per poc (efecte esglaó); un disseny gradual
        seria més just.
  \item Resposta oberta. Idea: substituir el salt per un tram decreixent continu, p.\,ex.
        $B(x)=600-c\,(x-\text{llindar})$ entre el llindar i el punt on arriba a $0$,
        triant $c$ perquè enllaci a tots dos extrems.
\end{sollist}
\end{solucions}
